\begin{myexem}
  \label{exem:run2}
  Let us reconsider our running example; see Example~\ref{exem:run1}.
  The optimal solution of Program~\ref{prog:primal}
  is represented by
  Figure~\ref{fig:primal} for $2d = 2$, 4, 6, 8, 10 and 12.
  \begin{figure}[!ht]
    \centering
    \tabulinesep=0.5mm
    \begin{tabu}{{X[0.85,c,m]X[c,m]X[c,m]X[c,m]X[c,m]}}
      & Node 1 & Node 2 & Node 3 & Node 4\\
      \PrimalFig{2}{1}
      \PrimalFig{4}{2}
      \PrimalFig{6}{3}
      \PrimalFig{8}{4}
      \PrimalFig{10}{5}
      \PrimalFig{12}{6}
    \end{tabu}
    \caption{
    Representation of the solutions to Program~\ref{prog:primal}
    with different values of $d$ for the running example.
    The blue curve represents the boundary of the
    1-sublevel set of the optimal
    solution $p_v$ at each node $v \in \Nodes$.
    The dashed curve is the boundary of the unit circle.
    Observe that some sets are not convex.}
    \label{fig:primal}
  \end{figure}
\end{myexem}
